% Generated by task tex-groundtruth from dossiers/gt-flock-ring.md (section 3, the checks),
% with the corrections of dossiers/verify-gt-flock-ring.md.
\subsection{The BLAKE2s validity phase (Flock) and ring switching: the checks}\label{sec:gen-flock-ring-checks}

\begin{sloppypar}
Short forms in this section; the number after the colon is the line, at leanVM's pin
\code{a386121f}.
\begin{itemize}
\item Specification, under \file{doc/leanvm/body/}: \code{c-flock:} is
  \file{c-flock-protocol.tex}, \code{a-ring:} \file{a-ring-switching.tex}, \code{03:}
  \file{03-proving-primitives.tex}, \code{06:} \file{06-bus-interactions.tex}, \code{08:}
  \file{08-end-to-end-protocol.tex}.
\item Rust, under \file{crates/flock/src/}: \code{hash.rs}, \code{zerocheck.rs},
  \code{lincheck.rs}; under \file{crates/flock/src/zerocheck/}: \code{zerocheck/multilinear.rs},
  \code{univariate_skip_optimized.rs}; under \file{crates/lean_vm/src/}: \code{cpu/mod.rs};
  under \file{crates/pcs/src/}: \code{stack_open.rs}; \code{transcript.rs} is
  \file{crates/fiat_shamir/src/transcript.rs}.
\item Python: \code{verifier.py} is \file{python-verifier/verifier.py}.
\end{itemize}
\end{sloppypar}

The phase has one equality check and a set of derivations that replace checks. A Lean verifier
that receives what the deployed one derives must check it; each row says what a verifier without
it would accept. The step numbers are those of \cref{tab:gen-flock-ring-messages}.

\begin{landscape}
{\small
\setlength{\tabcolsep}{4pt}
\begin{longtable}{@{}>{\raggedright\arraybackslash}p{0.15\linewidth}>{\raggedright\arraybackslash}p{0.13\linewidth}>{\raggedright\arraybackslash}p{0.15\linewidth}>{\raggedright\arraybackslash}p{0.11\linewidth}>{\raggedright\arraybackslash}p{0.18\linewidth}>{\raggedright\arraybackslash}p{0.22\linewidth}@{}}
\caption{The checks of the Flock phase and of ring switching, and the derivations that replace
checks.}\label{tab:gen-flock-ring-checks}\\
\toprule
Check & Specification & Rust & Python & Protects & Accepted without it \\
\midrule
\endfirsthead
\toprule
Check & Specification & Rust & Python & Protects & Accepted without it \\
\midrule
\endhead
\bottomrule
\endfoot
\textbf{Lincheck terminal identity} (step 11)
  & \code{c-flock:126-133}; \code{:270} \enquote{it is the only thing that checks anything}
  & \code{lincheck.rs:1252-1256}
  & \code{verifier.py:1176}
  & Everything upstream: $v_a, v_b, v_c$ are the values of $a, b, z$ at the point; the constant
  position; the link between the $s_i$ and the matrices
  & Any $\qflock$. The prover commits limb slots with a wrong digest and zeros elsewhere, sends
  arbitrary messages at steps 3, 6, 7, 9 and the \textbf{true} $s_i$; ring switching and the
  opening accept, since the $s_i$ are true of $q$ \\ \addlinespace[2pt]
\textbf{Zeros assumed on $H$} (step 5): the verifier interpolates with 64 zeros, it does not
  receive them
  & \code{c-flock:62}
  & \code{zerocheck/}\allowbreak\code{multilinear.rs:133-137}
  & \code{verifier.py:1144}
  & This is where $a·b + c = 0$ is imposed
  & If the 64 values on $H$ were read from the prover: the prover sends the true $P$ of a witness
  violating the R1CS; every later step is honest and passes \\ \addlinespace[2pt]
\textbf{$c_0$ derived} in each zerocheck round (step 6), equivalently the round check $c_0 +
  r_i(c_1 + c_2) = \mathit{claim}$
  & \code{03:110}, \code{c-flock:69}
  & \code{transcript.rs:296-302}
  & \code{verifier.py:411-413}
  & The chain from $v_P$ to $R_{\mathrm{zc}}$
  & With $c_0$ free: the prover ignores $v_P$, runs an honest sumcheck of the true (nonzero) sum,
  ends at true $v_a, v_b, v_c$ \\ \addlinespace[2pt]
\textbf{$v_c$ derived} (step 7), equivalently the check $R_{\mathrm{zc}} = v_a v_b + v_c$
  & \code{c-flock:70-73}
  & \code{zerocheck.rs:423}
  & \code{verifier.py:1149}
  & The link between the zerocheck and the lincheck
  & With $v_c$ sent and unchecked: true $v_a, v_b, v_c$ of a violating witness pass the lincheck;
  the zerocheck constrains nothing \\ \addlinespace[2pt]
\textbf{The $α³$ term} in the target (step 8) and in the terminal identity (step 11)
  & \code{c-flock:23-28}, \code{:112-118}, \code{:131}
  & \code{lincheck.rs:1202, 1240}
  & \code{verifier.py:1162, 1174}
  & $z(512, t) = 1$ for every compression
  & The all-zero block: every row of the R1CS is homogeneous, so $z_t = 0$ satisfies it
  (\code{hash.rs:1176-1183} \enquote{homogeneous rows accept zero without the pin}). The limb
  slots are then zero: the statement \enquote{the compression of the zero input under the zero
  chaining value is zero} is accepted \\ \addlinespace[2pt]
\textbf{$c_1$ derived} in each lincheck round (step 9), equivalently $c_1 + c_2 = \mathit{claim}$
  & \code{03:110}
  & \code{transcript.rs:298-299}
  & \code{verifier.py:409-410}
  & The chain from the target to $R_{\mathrm{lc}}$
  & With $c_1$ free: the prover runs the sumcheck of the true inner product of a wrong witness;
  the target is never compared \\ \addlinespace[2pt]
\textbf{$T$ computed by the verifier} from the $s_i$ (step 13)
  & \code{a-ring:45}
  & \code{stack_open.rs:522-524}
  & \code{verifier.py:1346}
  & The $s_i$ are the slices of the committed $\qflock$
  & With $T$ sent: the terminal identity is one linear equation in 64 unknowns; the prover picks
  $s$ satisfying it for any witness and sends the true $⟨W, q⟩$ \\ \addlinespace[2pt]
\textbf{The weighted claim enters the batch} (step 14)
  & \code{08:98-99}
  & \code{stack_open.rs:522-524}
  & \code{verifier.py:1413}
  & the same
  & the same \\ \addlinespace[2pt]
\textbf{The limb claims enter the batch} as claims on $\qflock$
  & \code{08:77, 97}
  & \code{cpu/mod.rs:798-806}, \code{stack_open.rs:525-527}
  & \code{verifier.py:892, 1413}
  & The limbs the bus used are the words Flock proved
  & A $\qflock$ of valid but unrelated compressions, and limb values chosen to balance the memory
  bus: a forged digest in memory \\ \addlinespace[2pt]
\textbf{The size floor} $τ_{\mathrm{BLAKE2S}} ≥ 3$ and the cap $≤ 32$
  & cap only, \code{06:80}; \textbf{the floor is absent}
  & \code{cpu/mod.rs:161-166}
  & \code{verifier.py:860-861}
  & The layout: $\qflock$ has at least $2^3$ blocks (\code{hash.rs:283-286}), and the limb columns
  share its instance cube
  & Nothing unsound is known: the mathematics needs no floor (the zerocheck needs
  $k_{\mathrm{bool}} ≥ 13$, \code{zerocheck.rs:354}, always true). A Lean verifier without it
  accepts size announcements the deployed ones reject. (The citation check adds: the Rust names
  the floor \enquote{the lincheck floor of \code{n_blocks_log ≥ 3} (\code{n_outer ≥ 8})},
  \code{hash.rs:281-282}; the lincheck verifier has no such check, \code{lincheck.rs:1151-1176}
  comparing lengths only, so the floor is a constraint of the prover's layout; what in the
  lincheck prover needs \code{n_outer ≥ 8} was not checked.) \\ \addlinespace[2pt]
\textbf{The stream is consumed}
  & \code{08:100} (implicit)
  & \code{cpu/mod.rs:769}
  & \code{verifier.py:1414}
  & No trailing data
  & Malleable proofs \\ \addlinespace[2pt]
Shape checks inside the reduction
  & none
  & \code{lincheck.rs:1151-1176}, \code{zerocheck.rs:354-356}
  & none
  & Internal consistency; they cannot fail on the deployed path (\code{hash.rs:990-1008} passes
  constants)
  & nothing \\
\end{longtable}
}
\end{landscape}

\begin{sloppypar}
Not a run-time check, and load-bearing: \textbf{the seven fixed coordinates have
$\F_2$-independent equality weights} (\code{03:95-100}, used at \code{c-flock:278}). The Rust
asserts it in a unit test (\code{univariate_skip_optimized.rs:920-960}); the probe of the
dossier's section 9.3 recomputes it with the Python field: rank 128 (the citation check re-ran the
probe and obtained the same output).
\end{sloppypar}
